\begin{abstract}
\noindent
People who use Pakistan Sign Language (PSL) and people who do not know it often have no shared language. SignBridge PK is a first step towards a tool that reads a single photograph of a hand and returns the most likely letter. We have no PSL photographs yet, so every model in this report is trained and tested on the American Sign Language (ASL) alphabet; the results describe ASL, not PSL. We built two kinds of recogniser for the 26 letters A to Z: a convolutional network (MobileNetV2) that reads the image, and support vector machine, random forest and neural network classifiers that read 21 hand landmarks found by MediaPipe. The data was split by groups of look-alike images, so that near-identical frames do not appear in both training and test. On the 1{,}560 test images the CNN reaches 0.9795 accuracy and 0.9795 macro-F1. The landmark models can answer only the 1{,}090 test images (69.9\%) in which a hand was detected; on those images the CNN scores 0.9780 and the chosen SVM 0.9688. These scores are optimistic. Many test images have a near-identical training image of the same letter (70.0\% within 8 bits of an image hash), the data has no signer labels, and we made no claim about fairness across people. The scores are not an estimate of accuracy on new signers or on PSL. The next step is to collect PSL photographs with signer labels and test on people not seen in training.
\end{abstract}
